Identificar y nombrar la superficie definida por cada función $z = f(x,y)$, describiendo su forma geométrica (paraboloide, silla de montar, plano, cono, cilindro, etc.) y señalando características clave como vértice, eje de simetría o punto de silla:
\begin{enumerate}
    \item $f(x,y) = 3x - 2y + 5$
    \item $f(x,y) = x^2 + 4y^2$
    \item $f(x,y) = \sqrt{x^2 + y^2}$
    \item $f(x,y) = 4 - x^2 - y^2$
\end{enumerate}